% GENERATED FILE. Edit canonical lessons, then run npm run generate:books.
% canonical-source-hash: fnv1a64:039f90850f96f17a
% canonical-lessons: ML-C29-suprabhaatham

\chapter{Good Morning}
\label{ch:ml-good-morning}
\hlchaptermodality{29}

\begin{chapteropening}
\textbf{By the end of this chapter:} \emph{I can greet someone in the morning in Malayalam with suprabhātaṁ.}

Say suprabhātaṁ, split su- from prabhātaṁ, and show it is NOT built on śubha the way good night is.
\end{chapteropening}

\section[\ml{സുപ്രഭാതം}]{\ml{സുപ്രഭാതം} (suprabhātaṁ) — Malayalam's real \textquotedblleft{}good morning,\textquotedblright{} a familiar word wearing a new hat}

\label{lesson:ML-C29-suprabhaatham}



\begin{quote}
You might expect Malayalam's \textquotedblleft{}good morning\textquotedblright{} to pair ML-C25's \textquotedblleft{}auspicious\textquotedblright{} word with ML-C26's morning word. It doesn't — Malayalam reaches for a completely different Sanskrit \textquotedblleft{}good,\textquotedblright{} and for a word you've technically already half-met.
\end{quote}

\subsection*{What to know first}
\textbf{\ml{സുപ്രഭാതം}} (\textbf{suprabhātaṁ}) — \textquotedblleft{}\textbf{good morning}\textquotedblright{} — is Malayalam's actual, most common morning greeting: two independently-fetched sources confirm it: one states it's used in \textbf{both formal and informal} contexts, the second separately calls it \textquotedblleft{}the most general greeting\textquotedblright{} for good morning — genuine corroboration, not a single source's claim taken on faith. Here's the honest surprise: it is \textbf{not} built on \textbf{\ml{ശുഭ}} (ML-C25's \textquotedblleft{}auspicious,\textquotedblright{} the word behind Malayalam's own \textquotedblleft{}good night\textquotedblright{}) at all. Instead, it's built from the Sanskrit \textbf{prefix su-} (\textquotedblleft{}good\textquotedblright{}) fused directly onto \textbf{\ml{പ്രഭാതം}} — already met, Chapter 26, as an alternative Malayalam word for \textquotedblleft{}morning.\textquotedblright{} Two genuinely different Sanskrit \textquotedblleft{}good\textquotedblright{} morphemes: \textbf{\ml{ശുഭ}} is a standalone adjective; \textbf{su-} here is a bound prefix that only ever attaches to another word.

\begin{culture}
Here's the honest cross-language twist: \textbf{\ml{സുപ്രഭാതം}} is the \textbf{exact same Sanskrit tatsama word} as Telugu's own \textbf{\te{సుప్రభాతం}} — same su- prefix, same \te{ప్రభాత}/prabhāta root. But the two languages hand it \textbf{opposite jobs}. In Malayalam, it's the \textbf{primary}, most commonly used morning greeting. In Telugu, it's only a \textbf{secondary} alternative — Telugu's primary greeting, \textbf{\te{శుభోదయం}}, is built on a completely different Sanskrit root (\emph{udaya}, \textquotedblleft{}rise,\textquotedblright{} not \emph{prabhāta}). Same borrowed word, opposite roles across two closely related languages.
\end{culture}

\subsection*{Your turn}
Say these aloud:

\begin{itemize}
\raggedright
  \item \textquotedblleft{}suprabhātaṁ\textquotedblright{} — \textquotedblleft{}good morning,\textquotedblright{} su- (\textquotedblleft{}good\textquotedblright{}) + prabhāta (\textquotedblleft{}shine, dawn,\textquotedblright{} ML-C26)
  \item the honest surprise — NOT built on śubha; a different Sanskrit \textquotedblleft{}good\textquotedblright{} morpheme entirely
  \item the cross-language inversion — Malayalam's primary greeting is Telugu's secondary alternative
\end{itemize}

\begin{tcolorbox}[breakable,colback=teal!4,colframe=teal!35!black,title={Before you move on}]
Is \ml{സുപ്രഭാതം} built on \ml{ശുഭ}, ML-C25's \textquotedblleft{}auspicious\textquotedblright{} word? (\textbf{No} — it's built on the Sanskrit prefix su-, a different \textquotedblleft{}good\textquotedblright{} morpheme, fused onto \ml{പ്രഭാതം}, ML-C26's alternative morning word.) Is \ml{സുപ്രഭാതം} the primary morning greeting in both Malayalam and Telugu? (\textbf{No} — it's PRIMARY in Malayalam but only a SECONDARY alternative in Telugu, where \te{శుభోదయం} — built on a different root — is primary.) Is the \textquotedblleft{}used in both formal and informal contexts\textquotedblright{} claim as shaky as Telugu's own morning-greeting claim? (\textbf{No} — it's confirmed by TWO independently-fetched sources, a genuinely stronger evidentiary footing than resting on just one.)
\end{tcolorbox}
